% concatenated from mainloglog.tex, canonicalheader.tex
\documentclass[11pt]{article}
\usepackage[margin=1in]{geometry}
\usepackage{mathpazo} % Palatino font with math support
\usepackage{amsmath,amssymb}
\usepackage{amsthm}
\usepackage{booktabs}
\usepackage{tikz}
\usetikzlibrary{calc,positioning}
\usepackage{xcolor}

% Visible referee/editorial comments added in blue; the original text below is unchanged.
\newenvironment{bluecomment}
  {\par\medskip\begingroup\color{blue}\small\noindent}
  {\par\endgroup\medskip}



\title{\bfseries A Log-Log Saving for Matrix-Algebra Length and Terseness}
\author{ Florian Ito Sprung\footnote{Current address: IHES, Le Bois-Marie, 35 route de Chartres  CS 40001,  91893 Bures-sur-Yvette, France\\ E-mail: ian.sprung@gmail.com}}
\date{}

% -----------------------------------------------------------------------------
% ALT TITLE 1: \title{\bfseries A Log-Log Saving in Matrix-Algebra Length}
% ALT TITLE 2: \title{\bfseries Stopping Rank Descent in Matrix Algebras}
% ALT TITLE 3: \title{\bfseries A Shorter Worst-Case Length Bound for Matrix Algebras}
% -----------------------------------------------------------------------------

\usepackage{amssymb}
\usepackage{amscd,amsfonts,amsmath,amssymb, bbm,dsfont,extarrows}
\usepackage{amsmath}
\usepackage{amsfonts}
\usepackage{amssymb}\usepackage{enumerate,epsf,fancyhdr,float,graphicx,tabularx}
\usepackage{latexsym,mathrsfs,multirow}
\usepackage{wasysym}
\usepackage{xypic}
\usepackage[all]{xy}\usepackage[OT2,T1]{fontenc}
%\usepackage[headings]{fullpage}
\usepackage[modulo,mathlines,displaymath, running]{lineno}
\usepackage{amsmath,mathrsfs,enumerate,enumitem,wasysym}
\usepackage{xypic,hhline}
\usepackage[all]{xy}
\usepackage[OT2,T1]{fontenc}
\usepackage[modulo,mathlines,displaymath]{lineno}
\usepackage{tikz,tikz-cd}
\usepackage{dashrule}
\usetikzlibrary{matrix}
%\addtolength{\textheight}{1.2cm}
%\addtolength{\voffset}{-0.6cm}
%\addtolength{\textwidth}{1cm}
%\addtolength{\hoffset}{-0.5cm}
\usepackage[modulo,mathlines,displaymath]{lineno}

% definitions specific to this author guide only
\newtheorem{theorem}{Theorem}[section]
\DeclareSymbolFont{cyrletters}{OT2}{wncyr}{m}{n}\DeclareMathSymbol{\Sha}{\mathalpha}{cyrletters}{"58}

%For margins
\newcommand{\mar}[1]{\marginpar{\raggedright\tiny #1}}

% little conveniences: one symbol
\newcommand{\bb}{{\beta}}
\newcommand{\e}{\varepsilon}
\renewcommand{\phi}{{\varphi}}
\newcommand{\Qp}{\mathbf{Q}_p}
\newcommand{\Zp}{\mathbf{Z}_p}
\newcommand{\Fp}{\mathbf{F}_p}
\newcommand{\ZZ}{\mathbf{Z}}
\newcommand{\Fpbar}{\overline{\mathbf{F}}_p}
\newcommand{\Ql}{\mathbf{Q}_l}
\newcommand{\Zl}{\mathbf{Z}_l}
\newcommand{\Fl}{\mathbf{F}_l}
\newcommand{\Flbar}{\overline{\mathbf{F}}_l}
\renewcommand{\geq}{\geqslant}
\renewcommand{\leq}{\leqslant}
\newcommand{\w}{{\widehat}}
\newcommand{\Ehat}{{\widehat{E}}}


% little conveniences: more than one sy
\newcommand{\smat}[1]{\left( \begin{smallmatrix} #1 \end{smallmatrix} \right)}
\newcommand{\mat}[1]{\left[ \begin{smallmatrix} #1 \end{smallmatrix} \right]}
\newcommand{\links}{\left(\begin{array}{cc}}
\newcommand{\rechts}{\end{array}\right)}
\newcommand{\bai}{\left[\begin{array}{cc}}
\newcommand{\dai}{\end{array}\right]}
\newcommand{\hidari}{\left(\begin{array}{c}}
\newcommand{\migi}{\end{array}\right)}
\newcommand{\lord}{\Z_p[\Delta][X]}
\newcommand{\congruent}{\equiv}
\newcommand{\roots}{\smat{-1 & -1\\  \beta & \alpha}}

% Bourbaki fonts
\newcommand{\A}{\mathbb{A}}
\newcommand{\Ad}{\mathbb{A}}
\newcommand{\C}{\mathbb{C}}
\newcommand{\DD}{\mathbb{D}}
\newcommand{\Id}{\mathbb{I}}
\newcommand{\M}{{\mathbb M}}
\newcommand{\MM}{\mathbb{M}}
\newcommand{\N}{\mathbb{N}}
\newcommand{\NN}{{\mathbb N}}
\newcommand{\Q}{\mathbb{Q}}
\newcommand{\Proj}{\mathbb{P}}
\newcommand{\R}{\mathbb{R}}
\newcommand{\Z}{\mathbb{Z}}

% Gothic fonts
\newcommand{\fg}{{\mathfrak g}}
\newcommand{\fh}{{\mathfrak h}}
\newcommand{\fM}{{\mathfrak M}}
\newcommand{\gA}{{\mathfrak A}}
\newcommand{\gp}{{\mathfrak p}}
\newcommand{\gq}{{\mathfrak q}}
\newcommand{\gM}{{\mathfrak M}}
\newcommand{\ga}{{\mathfrak a}}
\newcommand{\gb}{{\mathfrak b}}
\newcommand{\gc}{{\mathfrak c}}
\newcommand{\dg}{{\mathfrak g}}
\newcommand{\gf}{{\mathfrak f}}
\newcommand{\gh}{{\mathfrak h}}
\newcommand{\gH}{{\mathfrak H}}
\newcommand{\gm}{{\mathfrak m}}
\newcommand{\gP}{{\mathfrak P}}
\newcommand{\gs}{{\mathfrak s}}
\newcommand{\h}{{\mathfrak h}}
\newcommand{\II}{{\mathfrak O}}
\newcommand{\mm}{\mathfrak{m}}
\newcommand{\pp}{\mathfrak{p}}

% bold fonts
\newcommand{\kk}{\mathbf{k}}
\newcommand{\bz}{\mathbf{z}}
\newcommand{\bH}{\mathbf{H}}

% mathcal characters
\newcommand{\aA}{{\mathcal A}}
\newcommand{\calA}{\mathcal{A}}
\newcommand{\calB}{\mathcal{B}}
\newcommand{\calC}{\mathcal{C}}
\newcommand{\calD}{\mathcal{D}}
\newcommand{\calE}{\mathcal{E}}
\newcommand{\calF}{\mathcal{F}}
\newcommand{\calG}{\mathcal{G}}
\newcommand{\calH}{\mathcal{H}}
\newcommand{\calI}{\mathcal{I}}
\newcommand{\calJ}{\mathcal{J}}
\newcommand{\calK}{\mathcal{K}}
\newcommand{\calL}{\mathcal{L}}
\newcommand{\calM}{\mathcal{M}}
\newcommand{\calN}{\mathcal{N}}
\newcommand{\calO}{\mathcal{O}}
\newcommand{\calP}{\mathcal{P}}
\newcommand{\calQ}{\mathcal{Q}}
\newcommand{\calR}{\mathcal{R}}
\newcommand{\calS}{\mathcal{S}}
\newcommand{\calT}{\mathcal{T}}
\newcommand{\calU}{\mathcal{U}}
\newcommand{\calV}{\mathcal{V}}
\newcommand{\calW}{\mathcal{W}}
\newcommand{\calX}{\mathcal{X}}
\newcommand{\calY}{\mathcal{Y}}
\newcommand{\calZ}{\mathcal{Z}}
\newcommand{\D}{{\mathcal O}}
\newcommand{\E}{{\mathcal E}}
\newcommand{\F}{{\mathcal F}_{ss}}
\newcommand{\G}{{\mathcal G}}
\renewcommand{\H}{\mathcal{H}}
%\renewcommand{\l}{\textgt{l}}
\newcommand{\K}{\mathcal{K}}
\newcommand{\OO}{\mathcal{O}}
\newcommand{\UU}{\mathcal{U}}
\newcommand{\X}{{\mathcal X}}
\newcommand{\tor}{{\mathrm{tor}}}

% Very curly fonts
\newcommand{\FF}{\mathscr{F}}
\newcommand{\GG}{\mathscr{G}}
\newcommand{\JJ}{\mathscr{J}}
\newcommand{\XX}{\mathscr{X}}

% limits
\newcommand{\projlimn}{\displaystyle \varprojlim_n}
\newcommand{\projlimm}{\displaystyle \varprojlim_m}
\newcommand{\projlimN}{\displaystyle \varprojlim_N}
\newcommand{\projlimnu}{\displaystyle \varprojlim_\nu}
\newcommand{\indlimn}{\displaystyle \varinjlim_n}
\newcommand{\indlimm}{\displaystyle \varinjlim_m}
\newcommand{\indlimN}{\displaystyle \varinjlim_N}
\newcommand{\indlimnu}{\displaystyle \varinjlim_\nu}

% shortcuts from papers
\newtheorem{hatexample}[theorem]{$\widehat{\bf{Example}}$}
\newcommand{\Aux}{\tilde{Y}}
\newcommand{\dn}{{\delta_n^i}}
\newcommand{\dI}{\delta_n^{i+1}}
\newcommand{\di}{\delta_n^{i-1}}
\newcommand{\Log}{\mathcal Log}
\newcommand{\Loghat}{\widehat{\Log}_{\alpha,\beta}(1+T)}
\newcommand{\Logarithm}{\Log_{\alpha,\beta}(1+T)}
\newcommand{\Xloc}{\X_{loc}}
\newcommand{\Xcol}{\X_{Col}}

% Poonen's shortcuts
\newcommand{\Aff}{\mathbb{A}}
\newcommand{\bbH}{\mathbb{H}}
\newcommand{\usim}{\sim_u}
 
\newcommand{\PP}{\mathbb{P}}
\newcommand{\Qbar}{\overline{\Q}}
\newcommand{\Zhat}{\hat{\Z}}
\newcommand{\Zbar}{\overline{\Z}}
\newcommand{\kbar}{\overline{k}}
\newcommand{\Kbar}{\overline{K}}
\newcommand{\ksep}{k^{\operatorname{sep}}}
\newcommand{\Fbar}{\overline{\F}}
\newcommand{\Adeles}{\mathbf{A}}
%\newcommand{\congruent}{\equiv}

% Poonen's shortcuts I'd probably never use
\newcommand{\leftexp}[2]{{\vphantom{#2}}^{#1}{#2}}
\newcommand{\sigmaiota}{\leftexp{\sigma}{\iota}}
\newcommand{\sigmaphi}{\leftexp{\sigma}{\phi}}
\newcommand{\sigmatauphi}{\leftexp{\sigma\tau}{\phi}}
\newcommand{\tauphi}{\leftexp{\tau}{\phi}}

% Math operators
\DeclareMathOperator{\Ann}{Ann}
\DeclareMathOperator{\Br}{Br}
\DeclareMathOperator{\cd}{cd}
\DeclareMathOperator{\Cl}{Cl}
\DeclareMathOperator{\codim}{codim}
\DeclareMathOperator{\Cor}{Cor}
\DeclareMathOperator{\divv}{div}
\newcommand{\cond}{\operatorname{cond}}
\newcommand{\disc}{\operatorname{disc}}
\newcommand{\Disc}{\operatorname{Disc}}
\newcommand{\Reg}{\operatorname{Reg}}
\newcommand{\Stab}{\operatorname{Stab}}
\newcommand{\END}{{\EE}\!nd}
\DeclareMathOperator{\Ext}{Ext}
\newcommand{\EXT}{{\E}\!xt}
\DeclareMathOperator{\Frac}{Frac}
\DeclareMathOperator{\Gr}{Gr}
\newcommand{\HOM}{{\HH}\!om}
\newcommand{\ind}{\operatorname{ind}}
\DeclareMathOperator{\inv}{inv}
       %\def\proj{\mathop{{\mathrm proj}}\nolimits}
       %\def\Iw{\mathop{{\mathrm Iw}}\nolimits}
\DeclareMathOperator{\proj}{\mathrm proj}
\newcommand{\Iw}{\mathrm{Iw}}
\DeclareMathOperator{\lcm}{lcm}
\DeclareMathOperator{\Lie}{Lie}
\DeclareMathOperator{\nil}{nil}
\DeclareMathOperator{\Norm}{Norm}
\DeclareMathOperator{\Num}{Num}
\DeclareMathOperator{\PIC}{\textbf{Pic}}
\DeclareMathOperator{\Prob}{Prob}
\DeclareMathOperator{\PROJ}{\textbf{Proj}}
\DeclareMathOperator{\re}{Re}
\DeclareMathOperator{\rk}{rk}
\DeclareMathOperator{\scd}{scd}
\DeclareMathOperator{\SPEC}{\textbf{Spec}}
\DeclareMathOperator{\supp}{supp}
\newcommand{\Sym}{{\mathrm Sym}}
\DeclareMathOperator{\tr}{tr}
\DeclareMathOperator{\trdeg}{tr deg}

% Categories
\newcommand{\Ab}{\operatorname{\textbf{Ab}}}
\newcommand{\Groups}{\operatorname{\textbf{Groups}}}
\newcommand{\Schemes}{\operatorname{\textbf{Schemes}}}
\newcommand{\Sets}{\operatorname{\textbf{Sets}}}

% Operators
\newcommand{\Az}{\operatorname{Az}}
\newcommand{\ET}{\operatorname{\textbf{\'{E}t}}}
\newcommand{\fl}{\operatorname{f\textcompwordmark l}}
\newcommand{\op}{\operatorname{op}}
\newcommand{\perf}{\operatorname{perf}}
\newcommand{\red}{\operatorname{red}}
\newcommand{\sing}{\operatorname{sing}}
\newcommand{\smooth}{\operatorname{smooth}}
\newcommand{\tH}{\operatorname{th}}
\newcommand{\unr}{\operatorname{unr}}
\newcommand{\Zar}{\operatorname{Zar}}
\newcommand{\Cech}{\v{C}ech}
\newcommand{\GalQ}{{\Gal}(\Qbar/\Q)}
\newcommand{\HH}{\operatorname{H}}
\newcommand{\HHcech}{\check{\HH}}
\newcommand{\HHat}{\hat{\HH}}
\newcommand{\PGL}{\operatorname{PGL}}
\newcommand{\PSL}{\operatorname{PSL}}
\newcommand{\Spec}{\operatorname{Spec}}
\newcommand{\Mat}{\operatorname{Mat}}
\newcommand{\cha}{\operatorname{Char}}
\newcommand{\Tam}{\operatorname{Tam}}
\newcommand{\sha}{\operatorname{sha}}
\newcommand{\Gal}{\operatorname{Gal}}
\newcommand{\Jac}{\operatorname{Jac}}
\newcommand{\Div}{\operatorname{Div}}
\newcommand{\Pic}{\operatorname{Pic}}
\newcommand{\End}{\operatorname{End}}
\newcommand{\tors}{\operatorname{tors}}
\newcommand{\val}{\operatorname{val}}
\newcommand{\cor}{\operatorname{cor}}
\newcommand{\Ell}{{\mathcal L}}
\newcommand{\Frob}{\operatorname{Frob}}
\newcommand{\Hom}{\operatorname{Hom}}
\newcommand{\res}{\operatorname{res}}
\newcommand{\Res}{\operatorname{Res}}
\newcommand{\rank}{\operatorname{rank}}
\newcommand{\coker}{\operatorname{coker}}
\newcommand{\Aut}{\operatorname{Aut}}
\newcommand{\Ind}{\operatorname{Ind}}
\newcommand{\Char}{\operatorname{Char}}
\newcommand{\GL}{\operatorname{GL}}
\newcommand{\ord}{\operatorname{ord}}
\newcommand{\col}{\operatorname{Col}}
\newcommand{\loc}{\operatorname{loc}}
\newcommand{\Wan}{\operatorname{Wan}}
\newcommand{\un}{\operatorname{un}}
\newcommand{\ab}{\operatorname{ab}}
\newcommand{\im}{\operatorname{Im}}
\newcommand{\Sel}{\operatorname{Sel}}
\newcommand{\Tr}{\operatorname{Tr}}
\renewcommand{\O}{\mathcal{O}}
\newcommand{\sign}{\operatorname{sign}}
\newcommand{\SL}{\operatorname{SL}}
\newcommand{\Span}{\operatorname{Span}}

%Subscripts and superscripts
\newcommand{\an}{\mathrm{an}}
\newcommand{\et}{\textrm{\'{e}t}}
%\newcommand{\l}{\l{}}}

% human names for tex commands
\newcommand{\del}{\partial}
\newcommand{\directsum}{\oplus} % binary direct sum
\newcommand{\Directsum}{\bigoplus} % direct sum of a collection
\newcommand{\injects}{\hookrightarrow}
\newcommand{\intersect}{\cap} % binary intersection
\newcommand{\Intersection}{\bigcap} % intersection of a collection
\newcommand{\isom}{\simeq}
\newcommand{\notdiv}{\nmid}
\newcommand{\surjects}{\twoheadrightarrow}
\newcommand{\tensor}{\otimes} % binary tensor product
\newcommand{\Tensor}{\bigotimes} % tensor product of a collection
\newcommand{\union}{\cup} % binary union
\newcommand{\Union}{\bigcup} % union of a collection
\newcommand{\isomto}{\overset{\sim}{\rightarrow}}
\newcommand{\isomfrom}{\overset{\sim}{\leftarrow}}

%I don't understand what these are
\newcommand*{\code}[1]{\mdseries\texttt{#1}}
\newcommand*{\pkg}[1]{\mdseries\textsf{#1}}
\renewcommand{\topfraction}{0.99}
\renewcommand{\contentsname}{Contents\\{\footnotesize\normalfont(A table
of contents should normally not be included)}}

% names of statements
\newtheorem{acknowledgement}[theorem]{Acknowledgement}
\newtheorem{algorithm}[theorem]{Algorithm}
\newtheorem{application}[theorem]{Application}
\newtheorem{assumption}[theorem]{Assumption}
\newtheorem{auxiliary proposition}[theorem]{Auxiliary Proposition}
\newtheorem{axiom}[theorem]{Axiom}
\newtheorem{case}[theorem]{Case}
\newtheorem{choice}[theorem]{Choice}
\newtheorem{claim}[theorem]{Claim}
\newtheorem{conclusion}[theorem]{Conclusion}
\newtheorem{condition}[theorem]{Condition}
\newtheorem{conjecture}[theorem]{Conjecture}
\newtheorem{convention}[theorem]{Convention}
\newtheorem{corollary}[theorem]{Corollary}
\newtheorem{criterion}[theorem]{Criterion}
\newtheorem{definition}[theorem]{Definition}
\newtheorem{evidence}[theorem]{Evidence}
\newtheorem{example}[theorem]{Example}
\newtheorem{exercise}[theorem]{Exercise}
\newtheorem{false}[theorem]{False Proposition}
\newtheorem{lemma}[theorem]{Lemma}
\newtheorem{main conjecture}[theorem]{Main Conjecture}
\newtheorem{main theorem}[theorem]{Main Theorem}
\newtheorem{modesty proposition}[theorem]{Modesty Proposition}
\newtheorem{notation}[theorem]{Notation}
\newtheorem{observation}[theorem]{Observation}
\newtheorem{open problem}[theorem]{Open Problem}
\newtheorem{problem}[section]{Problem}
\newtheorem{propn}[theorem]{Proposition}
\newtheorem{proposition}[theorem]{Proposition}
\newtheorem{results}[theorem]{Results}
\newtheorem{remark}[theorem]{Remark}
\newtheorem{solultion}[theorem]{Solution}
\newtheorem{summary}[theorem]{Summary}

%\theoremstyle{lemmas with names}
\newtheorem{convergence lemma}[theorem]{Convergence Lemma}
\newtheorem{corrected lemma}[theorem]{Corrected Lemma}
\newtheorem{growth lemma}[theorem]{Growth Lemma}
\newtheorem{coefficient lemma}[theorem]{Integrality Lemma}
\newtheorem{interpolation lemma}[theorem]{Interpolation Lemma}
\newtheorem{kernel lemma}[theorem]{Kernel Lemma}
\newtheorem{limit lemma}[theorem]{Limit Lemma}
\newtheorem{tandem lemma}[theorem]{Modesty Lemma}
\newtheorem{zero-finding lemma}[theorem]{Zero-Finding Lemma}
\newtheorem{zerofindinglemma}[theorem]{Zero-Finding Lemma}


\newtheorem*{theoremA}{Theorem A}
\newtheorem*{theoremB}{Theorem B}

\numberwithin{equation}{section}

\begin{document}
\maketitle

% The paper makes one small but robust change in the \v{S}itov--Pappacena
% mechanism: stop the square-zero rank descent at the first rank at most
% log_2 n.  This saves the last log_2 log_2 n halvings of rank.
% PROOF ANSATZ COMMENT:

\begin{abstract}
We study how long words in a family of matrices must be before they linearly span the algebra generated by that family. Let $F$ be a field, $S\subseteq \Mat_n(F)$, and $FS^{\leq k}$ the $F$-linear span of all words in $S$ of length at most $k$. Set
$\ell(S):=\min\left\{k:FS^{\leq k}=F[S]\right\}.$
%the largest $\ell(S)$ over generating sets $S$, i.e.  
%$L(\Mat_n(F)):=\max\left\{\ell(S):S\subseteq \Mat_n(F), F[S]=\Mat_n(F)\right\}.$
\v{S}itov proved the estimate
$$ \ell(S)\leq 2n\log_2 n+4n-4. $$
We prove the log--log improvement
$$ \ell(S)\leq 2n\log_2 n-2n\log_2\log_2 n+5n $$
for every $n>1$ and every $S\subseteq \Mat_n(F)$.

%In fact, this is also a bound for $\mathcal L_n(F)=\max\left\{\ell(S):S\subseteq \Mat_n(F)\right\}$. 
A theorem of Specht gives a word-criterion for unitary similarity of complex $n \times n$ matrices.
The trace argument of Freedman–Gupta–Guralnick, as used by Pappacena, shows that our estimate can be used to bound the terseness $\tau(n)$, i.e. the shortest length of words needed in Specht’s theorem. Thus, for n > 1,
$$ \tau(n)\leq 4n\log_2 n-4n\log_2\log_2 n+10n+1,\qquad n>1. $$
%For the two-letter list of trace words, this replaces the order of magnitute $n^{2n}2^{O(n)}$ coming from \v{S}itov's cutoff by $(n/\log_2 n)^{2n}2^{O(n)}$.

\end{abstract}


\section{Introduction}

% One way to read the length problem is as a question about how many
% noncommutative products must be sampled before the linear span stops growing.
The length problem for matrix algebras asks how long one must multiply matrices until the resulting products span the \textit{algebra} generated by the given matrices as a \textit{vector space}.  More precisely, if $S\subseteq \Mat_n(F)$, let $F S^{\leq k}$ denote the $F$-linear span of all words in $S$ of length at most $k$, including the empty word (identity matrix). %"identity matrix" necessary? 
The length $\ell(S)$ is the least $k$ for which $FS^{\leq k}=F[S]$.  The  %worst-case
length of the full matrix algebra is $L(\Mat_n(F)):=\max_{F[S]=\Mat_n(F)}\ell(S).$ We can now state our main result.


\begin{theoremA}
	For every field $F$, every $n>1$, and every $S\subseteq\Mat_n(F)$, 
	$$	\ell(S)\leq 2n\log_2 n-2n\log_2\log_2 n+5n.$$
In particular, $L(\Mat_n(F))$ satisfies the same bound.
\end{theoremA}


% At this point one could introduce

The length problem goes back to work of Paz, who proved the quadratic estimate $\ell(S)\leq \left\lceil\frac{n^2+2}{3}\right\rceil$
in \cite[Theorem~1 and Remark~2]{Paz1984}.  Paz also conjectured the linear bound $\ell(S)\leq 2n-2;$
 this is Conjecture~2 of \cite{Shitov2019}, see also \cite[Conjecture~1.6]{GutermanLaffeyMarkovaSmigoc2018}.  The conjecture remains open in general.

The same question also appears in representation theory and invariant theory in the work of Freedman, Gupta, and Guralnick on \v{S}ir\v{s}ov's theorem and representations of semigroups, who ask for a bounding function $g(n)$ in \cite[Question~3.6]{FreedmanGuptaGuralnick1997} for the analogous spanning degree for semigroups of matrices. Their Corollary~2.8 specializes this to the unitary-similarity problem for a matrix and its adjoint as in Specht's theorem. Specht's theorem \cite[Satz 1]{Specht1940} says that two complex $n \times n$ matrices $A$ and $B$ are unitarily similar if and only if $\tr w(A,A^*)=\tr w(B,B^*)$ for all words $w$.  %Their Theorem~2.6 shows how the spanning degree gives us finite trace criteria for semigroup representations, and their Corollary~2.8 specializes this to the unitary-similarity problem for a matrix and its adjoint, as in Specht's theorem  Their Theorem~3.7 gives linear bounds in rank-one and simple-spectrum cases.

Inspired by this circle of ideas, Pappacena proved \cite{PappacenaThesis} a better bound for matrix length \cite[Corollary~3.2]{Pappacena1997}, giving the first subquadratic estimate
\begin{equation}
 \ell(S) < n\sqrt{\frac{2n^2}{n-1}+\frac14}+\frac n2-2 =\sqrt{2}\,n^{3/2}+O(n).\tag{*}
\end{equation}
Pappacena's crucial tool was a \textit{finishing lemma}, saying that if a short word span contains a matrix of a specified rank, that rank controls the number of additional letters in the words needed to span the full matrix algebra \cite[Theorem~4.1(a)]{Pappacena1997}. For another application of these length ideas in the context of affine semiprime algebras of Gelfand--Kirillov dimension one, see \cite[Theorem~5]{PappacenaSmallWald2003}. 

\v{S}itov \cite[Theorem~3 and Claim~14]{Shitov2019} sharpened the  estimate to
$$\ell(S)\leq 2n\log_2 n+4n-4 $$
by improving Pappacena's finishing lemma when the rank was one, and then crucially proving a \textit{descent} for square-zero matrices: Starting from a square-zero matrix, the descent repeatedly produces another square-zero matrix whose rank is at most half the previous rank, while keeping the `cost' for producing it low. The `cost' is the number of (extra) letters needed. \v{S}itov descends all the way to rank one, and then applies his improved finishing lemma. A careful reading of \cite[Proof of Theorem 3]{Shitov2019} then gives a slight improvement and results in the bound $$\ell(S)\leq 2n\log_2 n+2n-4. $$%explain more?

The main observation of this paper is that it would be more judicious to stop \v{S}itov's square-zero rank descent as soon as the rank falls below $\sqrt{2}\log_2 n.$ If one stops the descent at a rank between $2$ and $\sqrt{2}\log_2 n$, we obtain the bound by simply applying Pappacena's finishing lemma\footnote{in fact, the bound is slightly better}. However, if the descent falls directly to rank one, then this means the previous rank was larger than $\sqrt{2}\log_2 n$. \v{S}itov's descent idea, in simple terms, was to discover a smaller rank matrix by looking at a `longer' space (i.e. generated by longer words), and then bounding the \textit{rank-length}, i.e. the product of the rank of the desired matrix and the length of the space it lives in at each step. To make this useful at the final step, we use a telescoping argument which relies on the elementary fact that for $x\in(0,\frac{1}{2}]$, we have $ 1-x\leq -\frac{1}{2}\log_2 x$. This is the origin of the $\log$. % In the case of direct descent to rank one, t
The final bound  is of the form 
$$5n+2n\log_2 n-2n\log_2 (\text{previous rank}).$$ 
Since the final rank is 1, the final length is the same as the final rank-length, and thus can be bounded by the above formula with the comparatively large previous rank, which we recall was bigger than  $\sqrt{2}\log_2 n$. It is this observation that is responsible for the log-log saving. 
% In the application to Specht's theorem, the pair $\{A,A^*\}$ may not generate all of $\Mat_n(\C)$. With this in mind, we define $\calL_n(F):=\max_{S\subseteq \Mat_n(F)}\ell(S).$
%where $\ell(S)$ is the least word length needed to span the algebra $F[S]$ generated by $S$.  
%Note that $L(\Mat_n(F))\leq \calL_n(F).$
%We do not  need, that this inequality is  strict.  get the quantifiers right: the induction 

% The following is stated pair \{A,A^*\} may generate a proper subalgebra.


%Where does the term $-2n\log_2\log_2 n$ come from? \v{S}itov studies an important invariant in his descent, the \textit{rank-length}. Indeed, each halving of the square-zero rank `costs' about $2n$ in the rank-length product.  Descending from rank roughly $n$ all the way to rank $1$ therefore `costs' about $2n\log_2 n$. This is `too expensive' for our arguments, and the main idea is to `not pay' for the last $$\log_2\log_2 n $$ halvings, namely the halvings from rank $L$ down to rank $1$.  This results exactly in the saving $$2n\log_2\log_2 n. $$
%We chose to include the condition $n \geq 16=2^4$ only for convenience: since $L=\log_2 n$, we have $L\geq 4$ and $\log_2 L\geq 1$.  These are exactly the elementary inequalities needed in our arguments. CHECK FOR OVEREXPLANATIONS %The real mechanism is the logarithmic stopping rank $L=\log_2 n$.

The same bound has a direct consequence for the Specht--Pearcy trace invariants in the unitary-similarity problem.  Recall that Pearcy's theorem \cite[Theorem~1]{Pearcy1962} is an improvement on Specht's theorem,  saying that in Specht's criterion $\tr w(A,A^*)=\tr w(B,B^*),$ only words $w$ of length at most $2n^2$ have to be considered; see also Shapiro's discussion of Pearcy's bound \cite[p.~149]{Shapiro1991}. This gives way to the question how much the bound $2n^2$ can be improved: How \textit{terse} can the words be so that the conclusion of Specht's theorem still holds?
 Denoting by $A\usim B$ unitary similarity, we make the following definition:

\begin{definition}
The \textbf{terseness} $\tau(n)$ is the least integer $\geq0$ such that, for every $A,B\in\Mat_n(\C)$,  $A\usim B$ if and only if  $\tr w(A,A^*)=\tr w(B,B^*)$
for all words $w$ of length at most $\tau(n)$.
\end{definition}

In this terminology, Pearcy proved $\tau(n)\leq 2n^2$. If $g(n)$ satisfies
$\ell(S)\leq g(n)$ for every $S\subseteq\Mat_n(\C)$, then
Freedman--Gupta--Guralnick
\cite[Corollary~2.8]{FreedmanGuptaGuralnick1997} give us 
$\tau(n)\leq2g(n)+1$.  %Pappacena later observed \cite[Theorem~4.3]{Pappacena1997} that an upper bound for $\calL_n(\C)$ gives the same bound on $\tau(n)$ .
 Thus, Theorem A gives the following consequence.

\begin{theoremB}
For $n> 1$,
$$\tau(n)\leq4n\log_2 n-4n\log_2\log_2 n+10n+1.$$
Consequently, the list of two-letter words needed via Specht's criterion has size at most $\left(\frac{n}{\log_2 n}\right)^{4n}2^{O(n)}.$

%instead of $n^{2n}2^{O(n)}$ as in \v{S}itov.
\end{theoremB}

\begin{remark}[Exact powers versus length]
There is a related, but different, problem where one studies the
linear span $FS^k$ of words of exactly\footnote{If the identity is
adjoined to the alphabet, then this problem is equivalent to ours, because we can pad every word of length at most $k$ by copies of
$I$:
$F(S\union\{I\})^k=FS^{\leq k}.$} length $k$. 
Without adjoining $I$, however, this exact-power problem is 
different, since the spaces $FS^k$ may not form an increasing filtration.
%This is the point of view of the quantum Wielandt problem.  
Micha{\l}ek and \v{S}itov proved an $O(n^2\log n)$ bound for the corresponding exact-power stabilization problem \cite[Theorem~1]{MichalekShitov2019}. % The present paper concerns only the standard algebra-length filtration $FS^{\leq k}$.
A later preprint of
\v{S}itov proves that if
$FS^k=\Mat_n(F)$ holds for some $k$, then the least such $k$ is at most
$ n^2+2n-4 $
for $n>2$ \cite[Theorem~3]{ShitovGrowth2024}. 
\end{remark}

\section{Notation and known results}

Let $F$ be a field. We let all algebras be unital, and we let word spans include the empty word, which contributes the identity matrix.

\begin{definition}
Let $S\subseteq \Mat_n(F)$.  For $k\geq 0$, define the word span to be the $F$-vector space $$F S^{\leq k}:=\Span_F\{s_1s_2\cdots s_j:0\leq j\leq k,\ s_i\in S\}.$$
Let $F[S]$ denote the unital $F$-algebra generated by $S$.  Define the \textbf{length} $\ell(S)$ of $S$ to be the least $k$ such that $F S^{\leq k}=F[S].$
%We also define $\calL_n(F):=\max_{S\subseteq \Mat_n(F)}\ell(S).$
% Since $\Mat_n(F)$ is finite dimensional, the supremum is a maximum.  This notation is not meant to suggest that $L(\Mat_n(F))<\calL_n(F)$ occurs; no such inequality is needed.

% The distinction is purely involves
% arbitrary pairs \{A,A^*\}, which may generate a proper subalgebra.
\end{definition}

%The following theorem is rd irreduciblity criterion forreduction to the full matrix algebra.


Following\footnote{``Irreducible'' normally means that $S$ has no common
	nonzero proper invariant subspace.  The two definitions agree over algebraically closed fields.} \cite[1. Warm-Up]{Shitov2019}, we say that $S$ is \textit{irreducible} if $F[S]=\Mat_n(F)$. %A subalgebra $A\subseteq\Mat_n(F)$ \textit{acts irreducibly} on $F^n$ if there is no nonzero proper subspace of $F^n$ invariant under every element of $A$. 
A result\footnote{See Burnside's 1905 theorem on irreducible linear groups \cite{Burnside1905}. For the statement in matrix-algebra form; see the statement on the first page of Radjavi--Rosenthal \cite[Theorem 1.5.1]{RadjaviRosenthal2000}.} of Burnside says that if $F$ is algebraically closed and $S$ is not irreducible, then the elements of $S$ can be simultaneously put into upper block-triangular form.
 %field, such a subalgebra generates the full matrix algebra. % Burnside is the bridge from the irreducible case to the full matrix algebra.
% Reducible sets are handled by the block induction below.
%\begin{theorem}[Burnside]
%\label{thm:burnside}
%Let $F$ be algebraically closed and let $A\subseteq \Mat_n(F)$ be a unital $F$-subalgebra.  If $A$ acts irreducibly on $F^n$, then $A=\Mat_n(F)$.
%\end{theorem} %Burnside's theorem is useful for us because it says that the algebraically closed irreducible case is already handled by the full matrix algebra case. %As for the reducible cases, t
The following theorem bounds the length in terms of lengths coming from subblocks. This is \cite[Corollary~3]{Markova2005}, see also \cite[Lemma~4]{Shitov2019}. 

\begin{theorem}[Block reduction]
\label{thm:block-reduction}
Let $S\subseteq \Mat_n(F)$ be simultaneously block upper triangular, i.e. after a simultaneous change of basis, every $s\in S$ is of the form
$$ s=\smat{ s_p & * \\ 0 & s_q}$$
with $s_p\in \Mat_p(F)$, $s_q\in \Mat_q(F)$, and $p+q=n$. Put
$S_p:=\{s_p:s\in S\}, S_q:=\{s_q:s\in S\}.$

Then $\ell(S)\leq \ell(S_p)+\ell(S_q)+1.$
\end{theorem}

%Pappacena's general estimate supplies the finite base case.

%\begin{theorem}[Pappacena's general bound]
%\label{thm:pappacena-general}
%For every field $F$ and every $n\geq 2$, $$\calL_n(F) <n\sqrt{\frac{2n^2}{n-1}+\frac14}+\frac n2-2.$$
%In particular, $\calL_n(F)=O(n^{3/2})$.
%\end{theorem}

%This follows from Pappacena's general finite-dimensional algebra estimate \cite[Corollary~3.2]{Pappacena1997}, applied to subalgebras of $\Mat_n(F)$.

Pappacena's key estimate bounds $\ell(S)$ in terms of the rank of a matrix already found in a short word span. This is \cite[Theorem~4.1(a)]{Pappacena1997}.% \v{S}itov later improved the special \emph{rank-one finish}.

% The proof below really uses the full rank-r estimate for 2 <= r <= log_2 n.
% It uses only Pappacena's r=1 estimate in the overshoot case; the sharper
% \v{S}itov rank-one finish is recorded for context but not needed.
\begin{theorem}[Pappacena's finishing lemma]% and \v{S}itov's rank-one finishing lemma]
\label{thm:pappacena-rank}
Let $F$ be algebraically closed and $S\subseteq \Mat_n(F)$ be irreducible.  Suppose $F S^{\leq k}$ contains a matrix of rank $r>0$.  Then $\ell(S)\leq rn+n-r+k-1.
$
%In the special case $r=1$, we have for $k>2$,$$\ell(S)\leq 2n+k-4.$$
\end{theorem}

 % The rank-one improvement is \cite[Corollary~7]{Shitov2019}. In the proof of Theorem A below we use Pappacena's full rank estimate for ranks $2\leq r\leq \log_2 n$, and in the overshoot case its $r=1$ case, namely $\ell(S)\leq 2n+k-2$.  We do not need \v{S}itov's sharper rank-one refinement.

Pappacena also proved stronger conclusions when the generating set already contains a sufficiently convenient matrix, see \cite[Theorem~4.1(b),(c)]{Pappacena1997}. \v{S}itov addressed the excluded cases by showing the existence of square-zero matrices with lower and lower ranks that don't cost too much, i.e. with a controlled  number of letters. This is \textit{\v{S}itov's square-zero descent}:

\begin{theorem}[\v{S}itov's square-zero descent]
\label{thm:shitov-descent}
Let $F$ be algebraically closed and let $S\subseteq \Mat_n(F)$ be irreducible.
\begin{itemize} 
    \item There are nonzero integers $\lambda_0$ and $\rho_0$ so that $\lambda_0\rho_0\leq 2n$ and $F S^{\leq \lambda_0}$ contains a square-zero matrix of rank $\rho_0$.
     \item If $F S^{\leq \lambda_i}$ contains a square-zero matrix of rank $\rho_i\geq2$, then there are integers $\lambda_{i+1}$ and $\rho_{i+1}$ such that $F S^{\leq \lambda_{i+1}}$ contains a square-zero matrix of rank $\rho_{i+1}$, where $$1\leq \rho_{i+1}\leq \frac{\rho_i}{2} \text{ and }\lambda_{i+1}\leq\frac{\lambda_i\rho_i}{\rho_{i+1}}+\frac{4n(\rho_i-\rho_{i+1})}{\rho_i\rho_{i+1}}.$$
\end{itemize}
\end{theorem}

The initial square-zero matrix is \v{S}itov's Claim~11, and the descent step is \v{S}itov's Claim~14 \cite[Claims~11 and~14]{Shitov2019}.  We now spell out how the numerical estimate used later is extracted from this descent.

Apply Theorem~\ref{thm:shitov-descent} repeatedly, as long as the current rank is at least $2$.  Since the rank is at least halved at each step, this produces a finite sequence


$$(\lambda_0,\rho_0),\ (\lambda_1,\rho_1),\ldots,\ (\lambda_{t+1},\rho_{t+1}),$$
where $F S^{\leq \lambda_i}$ contains a square-zero matrix of rank $\rho_i$, and where $\rho_{t+1}=1$.  %For $i\geq 1$, the $i$-th descent step gives $\rho_i\leq \rho_{i-1}/2$.
% If$$\alpha_i:=\frac{\rho_i}{\rho_{i-1}},$$
%then multiplying the displayed inequality in Theorem~\ref{thm:shitov-descent} for the step from $i-1$ to $i$ by $\rho_i$ gives

Since $0<\frac{\rho_{i+1}}{\rho_{i}}\leq 1/2$, elementary calculus gives us
$ 1-\frac{\rho_{i+1}}{\rho_{i}}\leq -\frac{1}{2}\log_2 \frac{\rho_{i+1}}{\rho_{i}},$ so that 
\begin{equation}
	\label{eq:log-argument}
  \sum_{i=0}^j(1-\frac{\rho_{i+1}}{\rho_{i}})\leq\frac12\log_2\frac{\rho_0}{\rho_{j+1}}.
\end{equation}

 It is useful to multiply the word length by the current rank and set $\mu_i:=\lambda_i\rho_i.$ Thus $\mu_i$ is the \emph{rank-length} at the $i$-th stage of the descent.
 We  have  $ \mu_{i+1} \leq \mu_{i}+4n\left(1-\frac{\rho_{i+1}}{\rho_{i}}\right).$
Iterating this inequality and using $\mu_0\leq 2n$ gives
$ \mu_{j+1} \leq 2n+4n\sum_{i=0}^j\left(1-\frac{\rho_{i+1}}{\rho_{i}}\right).$ 
Applying equation (\ref{eq:log-argument}) to this, we obtain $%\begin{equation}
%\label{eq:mu-bound-weak}
  \mu_{j+1}\leq2n+2n\log_2\frac{\rho_0}{\rho_{j+1}}.$%\end{equation}
Since every square-zero $n\times n$ matrix has rank at most $n/2$, we have $\rho_0\leq n/2$. 
%Indeed, \eqref{eq:mu-bound-weak} and $\rho_0\leq n/2$ give
Thus, we have proved the following lemma:
\begin{lemma}\label{lem:mu-bound}The following inequality holds\footnote{It even holds at the zero index: $\mu_0\leq2n\leq2n\log_2\frac{n}{\rho_0}$}: $$\mu_{j+1} \leq 2n+2n\log_2\frac{n}{2\rho_{j+1}}=2n\log_2\frac{n}{\rho_{j+1}}.$$
\end{lemma}
 %Thus the slightly sharper form used in the proof is


Putting $j=t$ and remembering that $\rho_{t+1}=1$, Lemma \ref{lem:mu-bound} then gives us $ \lambda_{t+1}\leq 2n\log_2n.$ Combining this with a finishing lemma\footnote{\v{S}itov improves Pappacena's finishing lemma (Theorem \ref{thm:pappacena-rank}) in the rank one case from Pappacena's $\ell(S)\leq 2n+k-2$ to
$\ell(S)\leq 2n+k-4,$ cf. \cite[Corollary~7]{Shitov2019}.} culminates in \cite[Theorem~3]{Shitov2019}:

\begin{theorem}[\v{S}itov, slightly improved\footnote{\v{S}itov's published bound is the slightly worse $\ell(S)\leq 2n\log_2 n+4n-4$. %This is because \v{S}itov used a slightly weaker argument in \cite[Proof of Theorem 3]{Shitov2019} than (\ref{eq:log-argument}) and the line that follows, giving the slightly weaker bound $ \lambda_{t+1}\leq 2n\log_2n+2n$ NO WAIT IT'S ACTUALLY $ \lambda_{t+1}\leq 2n\log_2\rho_0+2n$, but n
Note that \v{S}itov ends \cite{Shitov2019} by writing `this paper does not show any effort on improving the $o(n \log n)$ part of the upper bound.' }]
\label{thm:shitov-global}
For every field $F$, every $n\geq2$, and every $S\subseteq\Mat_n(F)$, $$\ell(S)\leq2n\log_2n+2n-4. $$
%Equivalently, for every subset $S\subseteq \Mat_n(F)$, the algebra $F[S]$ is spanned by words in $S$ of length at most $2n\log_2 n+2n-4$.
\end{theorem}

%\begin{remark}
%    The astute reader may have noticed that the above arguments give the better bound $ \calL_n(F)\leq 2n\log_2 n+2n-4$
%\end{remark}
%CHEK IF THIS IS TRUE

We also record the following consequence.%Pappacena's comparison (relying on) with terseness.

%\begin{definition}
%For $n\geq 1$, define $\tau(n)$ to be the least integer $\tau$ such that, whenever $A,B\in \Mat_n(\C)$ satisfy $$\tr w(A,A^*)=\tr w(B,B^*)$$ for every word $w$ in two noncommuting letters $x,y$ with $|w|\leq \tau$, one has $A\usim B$.  This integer $\tau(n)$ is the \emph{terseness}.
%\end{definition}

\begin{theorem}\cite[Corollary~2.8]{FreedmanGuptaGuralnick1997}
\label{thm:specht-transfer}
Suppose $g(n)$ is a function such
that $\ell(S)\leq g(n)$ for every $S\subseteq\Mat_n(\C)$.  Then
$$ \tau(n)\leq2g(n)+1. $$ In words, any general upper bound for the length of all subsets of $\Mat_n(\C)$ gives a terseness bound of twice that bound plus one.
\end{theorem}


%The finite trace theorem behind it is due to Pearcy \cite[Theorem~1]{Pearcy1962}; the underlying infinite trace criterion is Specht's theorem \cite[Satz 1]{Specht1940}.

\section{The new result and proof}

We now prove the advertised improvement.  Give our bound a name by putting $$B(x):=2x\log_2 x-2x\log_2\log_2 x+5x\qquad(x\geq 2),$$ and set $B(1):=0$.  We need the following elementary fact to reduce the argument to the irreducible case, cf. the block reduction Theorem \ref{thm:block-reduction}.


\begin{lemma}
\label{lem:subadditive-B}
If $p,q\geq 1$ are integers and $p+q\geq 16$, then $ B(p)+B(q)+1\leq B(p+q).
$

\end{lemma}

Assuming this lemma, we now prove Theorem A.

\begin{theoremA}
Let $F$ be a field, $n>1$ be a positive integer, and $S\subseteq\Mat_n(F)$.  Then $$\ell(S)\leq2n\log_2 n-2n\log_2\log_2 n+5n.$$
In particular, $L(\Mat_n(F))\leq 2n\log_2 n-2n\log_2\log_2 n+5n$.
\end{theoremA}

\begin{proof}

% Strong induction is needed only to pass reducible sets to their diagonal blocks.
We induct on $n$, uniformly
over all fields and all matrix families, i.e. the $n$th induction step assumes that for every field $K$, every $m<n$, and every
$T\subseteq\Mat_m(K)$, one has $\ell(T)\leq B(m)$. 

% The irreducible algebraically closed case is where the rank descent 
%We prove the stronger all-subsets statement by induction on $n$.  Thus t
%The induction hypothesis says that . %for every field $K$ and every $m<n$.  
The cases $2\leq m<16$ follow from Pappacena's general bound $(*)$ from the introduction, see the table in the last section. So we assume $n\geq 16$. Let $S\subseteq \Mat_n(F)$. Let $\overline F$ be an algebraic
closure of $F$. We have $ \overline F\otimes_F F[S]\cong\overline F[S], \overline F\otimes_F FS^{\leq k}\cong\overline F S^{\leq k}.$
Faithful flatness then gives $FS^{\leq k}=F[S]$ if and only if $\overline F S^{\leq k}=\overline F[S]$, so
$\ell(S)$ is unchanged.  We consequently assume from now on that $F=\overline{F}$.

If $S$ is reducible, then by simultaneous triangularization there is a nontrivial block upper triangular form.  Let the diagonal block sizes be $p$ and $q$, with $p+q=n\geq 16$.  By Theorem~\ref{thm:block-reduction} and the induction hypothesis, $\ell(S) \leq B(p)+B(q)+1  \leq B(n)$ by Lemma~\ref{lem:subadditive-B}. 

 Thus we may assume that $F$ is algebraically closed and that $S$ is irreducible.  By Burnside's theorem, this means that $F[S]=\Mat_n(F)$, so the rank estimates and square-zero descent apply as in \cite{Shitov2019}. %We also assume that $n\geq 16$.

Put $$L:=\log_2 n, \text{ and } d(L):=\frac{B(n)}{n}=2L-2\log_2L+5.$$
Since $n\geq 16$, we have $L\geq 4$.  By the first step  in \v{S}itov's descent (Theorem~\ref{thm:shitov-descent}) there is a nonzero square-zero matrix of rank $\rho_0$ in $F S^{\leq \lambda_0}$ with $\lambda_0\rho_0\leq 2n.$

Pappacena's finishing lemma (Theorem \ref{thm:pappacena-rank}) gives
$$\ell(S)\leq \rho_0n+n+\lambda_0-\rho_0-1\leq n\left(\rho_0+1+\frac2\rho_0\right).
$$ When $\rho_0\leq \sqrt2 L$, this proves the theorem. Indeed, the function $x\mapsto x+1+2/x$ is convex, so on any closed interval,  it is
bounded by its endpoint values.  At $x=1$ it equals $4\leq d(L)$,
while at $x=\sqrt2 L$ it is at most
$\sqrt2 L+1+\sqrt2\leq d(L)$. Hence $\ell(S)\leq nd(L)=B(n)$ on the interval $[1,\sqrt2 L]$.

%Since $L\geq 4$, this is certainly at most $B(n)$. OPTIMIZE. 

We therefore assume $\rho_0>\sqrt2 L$. Starting from $(\lambda_0,\rho_0)$, apply the square-zero descent repeatedly, but only as long as the current rank is greater than $\sqrt2 L$.  %Note that the hypothesis $\rho \geq 2$ required by Theorem~\ref{thm:shitov-descent} is automatic since $L\geq 4$. 
Let $j$ be the first index such that $\rho_{j+1}\leq \sqrt2 L.$
There are two cases: either $2\leq \rho_{j+1}\leq \sqrt2 L$, or $\rho_{j+1}=1$.

\textbf{Case 1.} First suppose that $2\leq \rho_{j+1}\leq \sqrt2 L.$
By Lemma \ref{lem:mu-bound},
$$\lambda_{j+1}=\frac{\mu_{j+1}}{\rho_{j+1}}\leq\frac{2n\log_2(n/\rho_{j+1})}{\rho_{j+1}} =\frac{2n(L-\log_2\rho_{j+1})}{\rho_{j+1}}.$$
Using Pappacena's finishing lemma, Theorem~\ref{thm:pappacena-rank},
$$\ell(S)\leq n\rho_{j+1}+n+\lambda_{j+1}-\rho_{j+1}-1 \\\leq n\rho_{j+1}+n+\frac{2n(L-\log_2\rho_{j+1})}{\rho_{j+1}}.$$

Set $f_L(x):=x+1+\frac{2(L-\log_2 x)}{x}$. A direct differentiation gives $f_L''(x) =\frac{4(L-\log_2x)+6/\ln2}{x^3}>0$ for $2\leq x\leq \sqrt2 L$,
so $f_L$ is convex and it is enough to check the endpoints in $[2,\sqrt2 L]$.  At
$x=2,f_L(2)=L+2\leq d(L),$
because $L-2\log_2L+3\geq3$.  At $x=\sqrt2 L,$ $$f_L(\sqrt2 L)\leq \sqrt2 L+1+\frac{2L}{\sqrt2 L}=\sqrt2 L+1+\sqrt2\leq d(L).$$
Thus, $\ell(S)\leq nd(L)=B(n)$.

%A direct calculation gives L+3-f_L(r)=\frac{(r-2)(L-r)+2\log_2 r}{r}.

%I THINK $6n$ should be $7n$.

% f_L''(r)=(4(L-log_2 r)+6/log_e 2)/r^3>0, and then check the endpoints.
% The displayed identity above is shorter and gives the gap explicitly.
This is the desired estimate in the first case.

% This is the only dangerous case: the descent 
\textbf{Case 2.} It remains to treat the second case, the overshoot case $\rho_{j+1}=1.$

  By definition (minimality) of ${j}$, we have $\rho_j>\sqrt{2}L$.  The final descent step increases the rank-length $\mu$ by at most $4n$: indeed,
$\mu_{j+1}\leq\mu_{j}+4n\left(1-\frac{1}{\rho_j}\right)\leq\mu_{j}+4n.$
Using Lemma~\ref{lem:mu-bound} at the previous rank $\rho_j$,

$\mu_{j}\leq2n\log_2\frac{n}{\rho_j}\leq2n(L-\log_2 \sqrt2 L)=2nL-2n\log_2L-n,$
since $\rho_j>\sqrt2 L$. 
Since $\rho_{j+1}=1$, we have $\lambda_{j+1}=\mu_{j+1}$.  Hence
$\lambda_{j+1}\leq \mu_j+4n\leq 2nL-2n\log_2L-n+4n.$
The $r=1$ case of Pappacena's finishing lemma (Theorem \ref{thm:pappacena-rank}) now gives $$\ell(S)\leq2n+\lambda_{j+1}-2\leq2nL-2n\log_2 L+5n=B(n).$$

This proves the theorem.

\end{proof}

\begin{remark}
	The third term in our formula, $5n$, is optimal among integer multiples of $n$, if one stops the \v{S}itov descent at a constant multiple of $L$. We chose the multiple $\sqrt{2}L$ for convenience\footnote{Some exercises to the reader: Stopping at $aL$ gives the
coefficient $ 6-2\log_2 a$. The calculations become simpler when working simply with the trivial multiple $L$, where the third term becomes $6n$. %By allowing a not necessarily constant multiple, it seems that the third term can not be pushed below $4n+o(n)$ with the current methods.
}.
\end{remark}


% Theorem B is not a new invariant-theoretic argument; it is a direct transfer
% from the all-subsets matrix-length bound to Pearcy's finite trace criterion.
The consequence concerning the theorems of Specht and Pearcy is now immediate.

\begin{theoremB}
Let $n>1$.  Then
$$\tau(n)\leq4n\log_2 n-4n\log_2\log_2 n+10n+1.$$

\end{theoremB}


\begin{proof}
Apply Theorem~\ref{thm:specht-transfer} to Theorem A over $F=\C$.

\end{proof}

It remains to prove Lemma \ref{lem:subadditive-B}.

\begin{proof}[Proof of Lemma \ref{lem:subadditive-B}]
Write
$B(x)=x c(x)$ with
$$c(x):=2\log_2x-2\log_2\log_2x+5.$$
For $x\geq3,c'(x)=\frac{2}{x\ln2}\left(1-\frac1{\ln x}\right)>0,$ so $c$ is increasing, and therefore $c(n)\geq c(16)=9\text{ for }n\geq16.$
Moreover, $B''(x)=\frac{2}{x\ln2}\left(1-\frac1{\ln x}+\frac1{(\ln x)^2}\right)>0\text{for}x>1,$ so $B$ is strictly convex.

Put $n=p+q$.  If one of $p,q$ is $1$, say $p=1$, then $B(n)-B(n-1)=c(n)+(n-1)\bigl(c(n)-c(n-1)\bigr)\geq c(n)\geq9>1,$ which proves the claim.

Now suppose $p,q\geq2$.  The function $x\mapsto B(x)+B(n-x)$ is convex on $[2,n-2]$ (this comes down to differentiating $B(x)+B(n-x)$ twice), and hence its maximum is attained at an
endpoint.  Thus $B(p)+B(q)\leq B(2)+B(n-2).$
Since $B(2)=14$ and $B(n)-B(n-2)=2c(n)+(n-2)\bigl(c(n)-c(n-2)\bigr)\geq2c(n)\geq18>15,$ we obtain $B(p)+B(q)+1\leq B(2)+B(n-2)+1\leq B(n).$

\end{proof}

\section{A table comparing the bounds}

The following two tables compare the corresponding integer degree
cutoffs for {\small
$$ B(n):=2n\log_2 n-2n\log_2\log_2 n+5n,
B_{\mathrm{Pap}}(n):=n\sqrt{\frac{2n^2}{n-1}+\tfrac14}+\tfrac n2-2,
\text{ and }B_{\mathrm{Sh}}(n):=2n\log_2 n+2n-4. $$}
  

\begin{center}
\scriptsize
\setlength{\tabcolsep}{3pt}
\begin{tabular}{rrrr|rrrr|rrrr}
\toprule
$n$ & $\lfloor B\rfloor$ & $\lfloor B_{\mathrm{Pap}}\rfloor$ & $\lfloor B_{\mathrm{Sh}}\rfloor$ &
$n$ & $\lfloor B\rfloor$ & $\lfloor B_{\mathrm{Pap}}\rfloor$ & $\lfloor B_{\mathrm{Sh}}\rfloor$ &
$n$ & $\lfloor B\rfloor$ & $\lfloor B_{\mathrm{Pap}}\rfloor$ & $\lfloor B_{\mathrm{Sh}}\rfloor$ \\
\midrule
$2$ & $14$ & $\mathbf{4}$ & $\mathbf{4}$
&
$7$ & $53$ & $\mathbf{30}$ & $49$
&
$12$ & $101$ & $\mathbf{65}$ & $106$ \\

$3$ & $20$ & $\mathbf{8}$ & $11$
&
$8$ & $62$ & $\mathbf{36}$ & $60$
&
$13$ & $112$ & $\mathbf{73}$ & $118$ \\


$4$ & $28$ & $\mathbf{13}$ & $20$
&
$9$ & $72$ & $\mathbf{43}$ & $71$
&
$14$ & $122$ & $\mathbf{82}$ & $130$ \\

$5$ & $36$ & $\mathbf{18}$ & $29$
&
$10$ & $81$ & $\mathbf{50}$ & $82$
&
$15$ & $133$ & $\mathbf{90}$ & $143$ \\

$6$ & $44$ & $\mathbf{23}$ & $39$
&
$11$ & $91$ & $\mathbf{57}$ & $94$
&
$16$ & $144$ & $\mathbf{99}$ & $156$ \\
\bottomrule
\end{tabular}
\end{center}

The next table records some powers of two.  We also include $n=63$, since the integer cutoffs supplied by $B(n)$ and $B_{\mathrm{Pap}}(n)$ agree there, while $B(n)$ becomes strictly smaller at $n=64$.

\begin{center}
\scriptsize
\setlength{\tabcolsep}{3pt}
\begin{tabular}{rrrr|rrrr|rrrr}
\toprule
$n$ & $\lfloor B\rfloor$ & $\lfloor B_{\mathrm{Pap}}\rfloor$ & $\lfloor B_{\mathrm{Sh}}\rfloor$ &
$n$ & $\lfloor B\rfloor$ & $\lfloor B_{\mathrm{Pap}}\rfloor$  & $\lfloor B_{\mathrm{Sh}}\rfloor$ &
$n$ & $\lfloor B\rfloor$  & $\lfloor B_{\mathrm{Pap}}\rfloor$  & $\lfloor B_{\mathrm{Sh}}\rfloor$ \\
\midrule
$32$ & $331$ & $\mathbf{274}$ & $380$
&
$128$ & $\mathbf{1{,}713}$ & $2{,}119$ & $2{,}044$
&
$1{,}024$ & $\mathbf{18{,}796}$ & $46{,}876$ & $22{,}524$ \\

$63$ & $\mathbf{743}$ & $\mathbf{743}$ & $875$
&
$256$ & $\mathbf{3{,}840}$ & $5{,}931$ & $4{,}604$
&
$2{,}048$ & $\mathbf{41{,}126}$ & $132{,}130$ & $49{,}148$ \\

$64$ & $\mathbf{757}$ & $760$ & $892$
&
$512$ & $\mathbf{8{,}529}$ & $16{,}656$ & $10{,}236$
&
$4{,}096$ & $\mathbf{89{,}415}$ & $372{,}824$ & $106{,}492$ \\
\bottomrule
\end{tabular}
\end{center}


%$B(n)$ contains $\log_2\log_2 n$, while $P(n)$ contains the denominator $n-1$.  The displayed value $B_{\mathrm{sh}}(1)=-2$ is only the formal value of that expression.

\section*{Acknowledgments}
 The author thanks the Institut des Hautes Etudes Scientifiques for its hospitality and excellent working conditions. 
 
 %(to be inserted after acceptance of article.)


%The author thanks M for suggesting the term ``terseness,'' and Helene Shapiro for introducing the author to the theorems of Specht and Pearcy.  Thanks also go to the students in the author's fall 2025 Advanced Linear Algebra class for their enthusiasm and curiosity; Shapiro's \emph{Linear Algebra and Matrices: Topics for a Second Course} \cite{Shapiro2015} was the textbook for the course.  The author would also like to thank Robert Guralnick for pointing the author to the paper of Freedman--Gupta--Guralnick \cite{FreedmanGuptaGuralnick1997} and to Pappacena's thesis \cite{PappacenaThesis}.  The author thanks Yaroslav \v{S}itov for drawing his attention to the quadratic bound for the quantum Wielandt problem in \cite{ShitovGrowth2024}. Thanks also to Joseph Silverman, whoc helped improve the writing. Finally, the author thanks the Institut des Hautes Études Scientifiques for its hospitality during July 2026. 
The author was supported by Simons grant MPS-TSM-00008075.


\bibliographystyle{plain}
\bibliography{mainloglog}


\end{document}

Alternatie proof of lemma 1: \begin{proof}[Proof of Lemma\ref{lem:subadditive-B}]
Put $n=p+q$, and assume without loss of generality that $p\leq q$.  For $x>1$, write $$ B(x)=x c(x), c(x):=2\log_2 x-2\log_2\log_2 x+6 =2\log_2\!\left(\frac{x}{\log_2 x}\right)+6.
The function $c(x)$ is increasing for $x\geq 3$, since $x/\log_2 x$ is increasing there.  Also $c(16)=10$, and hence $c(n)\geq 10$ for $n\geq 16$. There are three cases. \begin{enumerate}
    
\item If $p=1$, then $$ B(n)-B(n-1) =n c(n)-(n-1)c(n-1) \geq c(n) \geq 10, $$ so $B(1)+B(n-1)+1\leq B(n)$.

\item If $p=2$, then $B(2)=16$, and $$ B(n)-B(n-2) =n c(n)-(n-2)c(n-2) \geq 2c(n)\geq 20 >B(2)+1. $$
Thus $B(2)+B(n-2)+1\leq B(n)$.

\item Finally assume $p\geq 3$.  Then $q\geq 3$ and $p\leq n/2$.  Since $c(x)$ is increasing for $x>3$, $$ B(n)-B(p)-B(q)=p\bigl(c(n)-c(p)\bigr)+q\bigl(c(n)-c(q)\bigr)\geq p\bigl(c(n)-c(p)\bigr)\geq p\bigl(c(n)-c(n/2)\bigr)>p\geq 3. $$

The last above inequality follows from
$$c(n)-c(n/2)=2\log_2\left(\frac{n}{L}\right)-2\log_2\left(\frac{n}{2(L-1)}\right)=2-2\log_2\left(\frac{L}{L-1}\right)=2\log_2\left(\frac{2(L-1)}{L}\right)=2\log_2\left(2-\frac{2}{L}\right)\geq 2\log_2\left(\frac32\right)>1,$$
  In the last inequality, we used $L=\log_2 n\geq 4$.
  
In particular $B(p)+B(q)+1\leq B(n)$.

\end{enumerate} 

\end{proof}
